\documentclass[a4paper]{article}

%% Language and font encodings
\usepackage[english]{babel}
\usepackage[utf8x]{inputenc}
\usepackage[T1]{fontenc}
\usepackage{ mathrsfs }
\usepackage{amsfonts}

%% Sets page size and margins
\usepackage[a4paper,top=3cm,bottom=2cm,left=3cm,right=3cm,marginparwidth=1.75cm]{geometry}

%% Useful packages
\usepackage{amsmath}
\usepackage{graphicx}
\usepackage[colorinlistoftodos]{todonotes}
\usepackage[colorlinks=true, allcolors=blue]{hyperref}
\title{Informal Notes on Rank of a Matrix}
\author{Harsha Vardhan Tetali}
\begin{document}
\maketitle
\author

\section{Introduction}

Let $A$ be a matrix of m rows and n columns with real entries ($A \in \mathbb{R}^{m \times n}$). We define two quantities called the \textbf{Row rank} and \textbf{Column Rank} of the matrix. 

\subsection*{Row Rank}

If each row of the matrix $A$ is looked as a vector in $\mathbb{R}^n$, then the row rank of the matrix is the number of linearly independent rows in the matrix.

We know from matrix vector multiplication that pre-multiplying a matrix with a row vector gives a row vector which is a linear combination of the rows of the matrix.

For example, let $A = \begin{bmatrix}
\boldsymbol{r}_1^{\top} \\ \boldsymbol{r}_2^{\top} \\ \vdots \\ \boldsymbol{r}_m^{\top} 
\end{bmatrix}$ and $\boldsymbol{y}^{\top} = \begin{bmatrix}
y_1 & y_2 & \cdots & y_m
\end{bmatrix}$ where for $i=1,2,\cdots,m$ we have $\boldsymbol{r}_i \in \mathbb{R}^{n}$.

We know that \begin{equation}
\boldsymbol{y}^{\top} A = y_1\boldsymbol{r}_1^{\top} + y_2\boldsymbol{r}_2^{\top} + \cdots + y_m\boldsymbol{r}_m^{\top} 
\end{equation}

Now let us take vectors $\boldsymbol{y}_1, \boldsymbol{y}_2, \cdots, \boldsymbol{y}_k \in \mathbb{R}^{m}$. Let us follow the following convention for each of them:

\begin{equation}
\boldsymbol{y}_i^{\top} = \begin{bmatrix}
y_{i1} & y_{i2} & \cdots & y_{im}
\end{bmatrix}
\end{equation}
for $i = 1,2,\cdots,k$.

It is easy to see that each $\boldsymbol{y}_i^{\top}A \in \mathbb{R}^{1 \times n}$. We can place all these row vectors as rows of a new matrix to create,

$\begin{bmatrix}
\boldsymbol{y}_1^{\top}A \\
\boldsymbol{y}_2^{\top}A \\
\vdots \\
\boldsymbol{y}_k^{\top}A
\end{bmatrix} = B$ where $B \in \mathbb{R}^{k \times n}$.

This entire process described above can be denoted in a simpler way:

\begin{equation}
\begin{bmatrix}
\boldsymbol{y}_1^{\top} \\
\boldsymbol{y}_2^{\top} \\
\vdots \\
\boldsymbol{y}_k^{\top}
\end{bmatrix}A = B
\end{equation}

Elaborating further (for lucidity) we can write,

\begin{equation}
\begin{bmatrix}
y_{11} & y_{12} & \cdots & y_{1m} \\
y_{21} & y_{22} & \cdots & y_{2m} \\
\vdots & \vdots & \ddots & \vdots \\
y_{k1} & y_{k2} & \cdots & y_{km}
\end{bmatrix}A = B
\end{equation}


Naming the first matrix in the left hand side above, as $Y \in \mathbb{R}^{k \times m}$, we can write.

\begin{equation}
YA = B
\end{equation}

In the above process, we observe:

\begin{enumerate}
\item A definition for the pre-multiplication of a matrix $Y$ with matrix $A$.
\item Each row of matrix $B (=YA)$ is a linear combination of rows of matrix A.
\end{enumerate}

From the second point above we can say that the row rank of matrix $B$ is utmost the rank of matrix $A$.

Therefore, to summarize:

$\text{row-rank}(B) \leq \text{row-rank}(A)$

or, \begin{equation}
\text{row-rank}(YA) \leq \text{row-rank}(A)
\end{equation}

Thus, pre-multiplication can either reduce or utmost maintain the rank of a matrix.

\subsection*{Column Rank}

If each column of the matrix $A$ is looked as a vector in $\mathbb{R}^m$, then the column rank of the matrix is the number of linearly independent columns in the matrix.

Assuming the same setup of matrix $A$ we now denote it with its columns: $A = \begin{bmatrix}
\boldsymbol{c}_1 & \boldsymbol{c}_2 & \cdots & \boldsymbol{c}_n 
\end{bmatrix}$.

We follow a similar exercise as in the previous subsection.

In this setup, now define a vector $\boldsymbol{x}= \begin{bmatrix}
x_1 \\
x_2 \\
\vdots \\
x_n
\end{bmatrix}$

Post-multiplying $A$ with $\boldsymbol{x}$, we get a linear combination of the columns of the matrix $A$.

\begin{equation}
A \boldsymbol{x} = x_1\boldsymbol{c}_1 + x_2\boldsymbol{c}_2 + \cdots + x_n\boldsymbol{c}_n 
\end{equation}

Now we assume $k$ vectors $\boldsymbol{x}_1, \boldsymbol{x}_2, \cdots, \boldsymbol{x}_k \in \mathbb{R}^n$ each represented as:

\begin{equation}
\boldsymbol{x}_i = \begin{bmatrix}
x_{1i}\\
x_{2i}\\
\vdots \\
x_{ni}
\end{bmatrix}
\end{equation}

Now we observe that each $A\boldsymbol{x}_i \in \mathbb{R}^{m}$ for all $i=1,2,\cdots,k$.

Now placing these above vectors in the columns of a new matrix we get,

\begin{equation} 
\begin{bmatrix}
A\boldsymbol{x}_1 & A\boldsymbol{x}_2 & \cdots & A\boldsymbol{x}_k 
\end{bmatrix} = C
\end{equation}

where $C \in \mathbb{R}^{m \times k}$. 
To represent the above in lucid way, we write,

\begin{equation} 
A\begin{bmatrix}
\boldsymbol{x}_1 & \boldsymbol{x}_2 & \cdots & \boldsymbol{x}_k 
\end{bmatrix} = C
\end{equation}

In an elaborate way we have,

\begin{equation}
A\begin{bmatrix}
x_{11} & x_{12} & \cdots & x_{1k} \\
x_{21} & x_{22} & \cdots & x_{2k} \\
\vdots & \vdots & \ddots & \vdots \\
x_{n1} & x_{n2} & \cdots & x_{nk}
\end{bmatrix} = C
\end{equation}

Naming the second matrix in the left hand side above, as $X \in \mathbb{R}^{n\times k}$, we can write:

\begin{equation}
AX = C
\end{equation}

In the above process we again examine two things:

\begin{enumerate}
\item A definition for the post-multiplication of a matrix $X$ with matrix $A$.
\item Each column of matrix $C(=AX)$ is a linear combination of columns of $A$.
\end{enumerate}
Again, from the second point above, we can say that the column rank of matrix $C$ is utmost the column rank of $A$.

Therefore, to summarize:

$\text{column-rank}(C) \leq \text{column-rank}(A)$

or, \begin{equation}
\text{column-rank}(AX) \leq \text{column-rank}(A)
\end{equation}

\section{Deficient Row Rank Matrices}

We call a matrix to have deficient row rank if the row-rank is less than the number of rows in the matrix. In such a case it is clear that the rows form a linearly dependent set. Considering the matrix $A \in \mathbb{R}^{m \times n}$, if the row-rank($A$) = r, say, is less than m (number of rows in $A$), i.e. $r < m$, then, it implies that $m-r$ of the rows can be represented as the linear combination of the $r$ linearly independent rows of $A$.

Let us follow the row convection as described in the subsection on row rank.

\begin{equation}
A = \begin{bmatrix}
\boldsymbol{r}_1^{\top} \\ \boldsymbol{r}_2^{\top} \\ \vdots \\ 
\boldsymbol{r}_{r}^{\top} \\
\boldsymbol{r}_{r+1}^{\top} \\ \vdots \\ 
\boldsymbol{r}_m^{\top} 
\end{bmatrix}
\end{equation}

Without loss of generality, let the first $r$ rows be linearly independent, i.e. $\{\boldsymbol{r}_{1}^{\top}, \boldsymbol{r}_{2}^{\top}, \cdots, \boldsymbol{r}_{r}^{\top}\}$ forms a linearly independent set.

Therefore the vectors $\boldsymbol{r}_{r+1}^{\top}, \boldsymbol{r}_{r+2}^{\top}, \cdots, \boldsymbol{r}_{m}^{\top}$ can be written as a linear combination of vectors from the above set.

Let,
\begin{equation}
\boldsymbol{r}_{i}^{\top} = a_{i,1}\boldsymbol{r}_{1}^{\top} + a_{i,2}\boldsymbol{r}_{2}^{\top} + \cdots + a_{i,r}\boldsymbol{r}_{r}^{\top}
\end{equation}

for $i = r+1, r+2, \cdots m$.

Thus we can write matrix $A$ as,

\begin{equation}
A = \begin{bmatrix}
1.\boldsymbol{r}_1^{\top}+ 0.\boldsymbol{r}_2^{\top} + \cdots + 0.\boldsymbol{r}_r^{\top}   \\ 0.\boldsymbol{r}_1^{\top}+ 1.\boldsymbol{r}_2^{\top} + \cdots + 0.\boldsymbol{r}_r^{\top}  \\ \vdots \\ 
0.\boldsymbol{r}_1^{\top}+ 0.\boldsymbol{r}_2^{\top} + \cdots + 1.\boldsymbol{r}_r^{\top} \\
a_{r+1,1}\boldsymbol{r}_{1}^{\top} + a_{r+1,2}\boldsymbol{r}_{2}^{\top} + \cdots + a_{r+1,r}\boldsymbol{r}_{r}^{\top}\\ 
a_{r+2,1}\boldsymbol{r}_{1}^{\top} + a_{r+2,2}\boldsymbol{r}_{2}^{\top} + \cdots + a_{r+2,r}\boldsymbol{r}_{r}^{\top}\\ \vdots \\ 
a_{m,1}\boldsymbol{r}_{1}^{\top} + a_{m,2}\boldsymbol{r}_{2}^{\top} + \cdots + a_{m,r}\boldsymbol{r}_{r}^{\top}
\end{bmatrix}
\end{equation}

Using intuitions obtained in the subsection of row-rank, we can factorize the above matrix as:

\begin{equation}
A = \begin{bmatrix}
1 & 0 & \cdots & 0   \\ 
0 & 1 & \cdots & 0   \\ 
\vdots & \vdots & \ddots & \vdots\\ 
0 & 0 & \cdots & 1 \\
a_{r+1,1} & a_{r+1,2} &\cdots & a_{r+1,r}\\ 
a_{r+2,1} &  a_{r+2,2} &\cdots & a_{r+2,r}\\ 
\vdots & \vdots & \ddots & \vdots\\ 
a_{m,1}& a_{m,2}& \cdots & a_{m,r}
\end{bmatrix} \begin{bmatrix}
\boldsymbol{r}_1^{\top} \\
\boldsymbol{r}_2^{\top} \\
\vdots \\
\boldsymbol{r}_r^{\top}
\end{bmatrix}
\end{equation}

Let \begin{equation}
S = \begin{bmatrix}
1 & 0 & \cdots & 0   \\ 
0 & 1 & \cdots & 0   \\ 
\vdots & \vdots & \ddots & \vdots\\ 
0 & 0 & \cdots & 1 \\
a_{r+1,1} & a_{r+1,2} &\cdots & a_{r+1,r}\\ 
a_{r+2,1} &  a_{r+2,2} &\cdots & a_{r+2,r}\\ 
\vdots & \vdots & \ddots & \vdots\\ 
a_{m,1}& a_{m,2}& \cdots & a_{m,r}
\end{bmatrix}
\end{equation}

and 
\begin{equation}
R = \begin{bmatrix}
\boldsymbol{r}_1^{\top} \\
\boldsymbol{r}_2^{\top} \\
\vdots \\
\boldsymbol{r}_r^{\top}
\end{bmatrix}
\end{equation}

Observe that each of the matrices $S$ and $R$ above have row-rank $r$ since both have $r$ independent rows.

Now let us see an interesting result. From the property of post multiplication we have:

\begin{equation}
\text{column-rank}(A) = \text{column-rank}(SR) \leq \text{column-rank}(S)
\end{equation}

Observe that the columns of $S$ are independent and are $r$ in number. 

Therefore,
\begin{equation}
\text{column-rank}(S) = r = \text{row-rank}(A)
\end{equation}

Therefore, we finally have,

\begin{equation}
\text{column-rank}(A) \leq \text{row-rank}(A)
\label{eq:fir}
\end{equation}

This result is true for all matrices, therefore for a matrix $A^{\top}$, we have,

\begin{equation}
\text{column-rank}(A^{\top}) \leq \text{row-rank}(A^{\top})
\end{equation}

It is now simple to see that,

\begin{equation}
\text{row-rank}(A) \leq \text{column-rank}(A)
\label{eq:sec}
\end{equation}

Thus, from \ref{eq:fir} and \ref{eq:sec} we have,

\begin{equation}
\text{row-rank}(A) = \text{column-rank}(A)
\end{equation}

Note that the above result to hold, the matrix need not be rank deficient, as in such a case, the values $a_{i,j}$ are all zeros.

Thus, we finally arrive a point where we see the two notions \textbf{row rank} and \textbf{column rank} of a matrix are the same number. From this, we define the rank of matrix which is equal to its row-rank and the column-rank.




\end{document}.


